\subsection{DEFINITION OF COXETER TRANSFORMATIONS}

An ordered chamber relative to W is a pair consisting of a chamber C determined by H and a bijection \(i \mapsto H_{i}\) from \(\{1, 2, \ldots, l\}\) onto the set of walls of C (cf.~§3, no. 9, Prop. 7).

DEFINITION 1. The Coxeter transformation determined by an ordered chamber \((\mathrm{C},(\mathrm{H}_i)_{1\leqslant i\leqslant l})\) is the element \(c = s_{\mathrm{H}_1}s_{\mathrm{H}_2}\dots s_{\mathrm{H}_l}\) of W.

PROPOSITION 1. All Coxeter transformations are conjugate in W.

Since W permutes the chambers determined by H transitively (§3, no. 2, Th. 1), we are reduced to proving the following: let \((\mathrm{C},(\mathrm{H}_{i})_{1\leqslant i\leqslant l})\) be an ordered chamber and \(\pi\in S_{l}\) ; then \(s_{H_{1}}s_{H_{2}}\ldots s_{H_{l}}\) and \(s_{H_{\pi(1)}}s_{H_{\pi(2)}}\ldots s_{H_{\pi(l)}}\) are conjugate in W. Taking into account §4, no. 8, Prop. 8, this will follow immediately from the following lemma:

Lemma 1. Let X be a finite forest, and \(x \mapsto g_{x}\) a map from X to a group \(\Gamma\) such that \(g_{x}\) and \(g_{y}\) are conjugate whenever x and y are not linked in X. Let T be the set of total orderings on X. For all \(\xi \in T\) , let \(p_{\xi}\) be the product in \(\Gamma\) of the sequence \((g_{x})_{x \in X}\) defined by \(\xi\) . Then the elements \(p_{\xi}\) are conjugate in \(\Gamma\) .

\begin{enumerate}
\def\labelenumi{\arabic{enumi})}
\item
  We proceed by induction on \(n = \operatorname{Card} X\) . The case \(n = 1\) is immediate, so assume that \(n \geqslant 2\) . There exists in \(X\) a terminal vertex \(a\) (Chap. IV, Appendix, no. 3, Prop. 2). Let \(b \in X - \{a\}\) be a vertex linked to \(a\) if one exists; if \(a\) is not linked to any vertex in \(X - \{a\}\) , let \(b\) in \(X - \{a\}\) be arbitrary. In all cases, \(g_a\) commutes with \(g_x\) for \(x \neq b\) . Let \(\eta \in T\) be such that \(a\) is the largest element of \(X\) and \(b\) the largest element of \(X - \{a\}\) ; we let \(\xi \in T\) and prove that \(p_\xi, p_\eta\) are conjugate.
\item
  Assume first that, for \(\xi\) , \(a\) is the largest element of X and \(b\) the largest element of \(\mathrm{X} - \{a\}\) . Let \(\mathrm{X}'\) be the full subgraph \(\mathrm{X} - \{a\}\) , which is a forest. Define a map \(x \mapsto g_x'\) from \(\mathrm{X}'\) to \(\Gamma\) by putting \(g_x' = g_x\) if \(x \neq b\) , \(g_b' = g_b g_a\) . Let \(\xi', \eta'\) be the restrictions of \(\xi, \eta\) to \(\mathrm{X}'\) . The induction hypothesis applies, so \(p_{\xi'}\) and \(p_{\eta'}\) are conjugate. But it is clear that \(p_{\xi'} = p_{\xi}\) , \(p_{\eta'} = p_{\eta}\) , proving the lemma in this case.
\item
  Assume that \(a\) is the largest element of \(X\) for \(\xi\) . Let \(X_1\) (resp. \(X_2\) ) be the set of elements of \(X - \{a\}\) strictly larger (resp. smaller) than \(b\) ; let \(\xi_i\) be the restriction of \(\xi\) to \(X_i\) . Then
\end{enumerate}

\[
p _ {\xi} = p _ {\xi_ {1}} g _ {b} p _ {\xi_ {2}} g _ {a} = p _ {\xi_ {1}} g _ {b} g _ {a} p _ {\xi_ {2}},
\]

and this element is conjugate to \(p_{\xi_2}p_{\xi_1}g_b g_a\) . We are thus reduced to case 2).

\begin{enumerate}
\def\labelenumi{\arabic{enumi})}
\setcounter{enumi}{3}
\tightlist
\item
  In the general case, let \(\mathrm{X}_3\) (resp. \(\mathrm{X}_4\) ) be the set of elements of \(\mathbf{X}\) strictly larger (resp. smaller) than \(a\) ; let \(\xi_i\) be the restriction of \(\xi\) to \(\mathrm{X}_i\) . Then \(p_{\xi} = p_{\xi_3}g_a p_{\xi_4}\) , and this element is conjugate to \(p_{\xi_4}p_{\xi_3}g_a\) . We are thus reduced to case 3).
\end{enumerate}

It follows from Prop. 1 that all the Coxeter transformations are of the same order \(h = h(\mathrm{W})\) . This number is called the Coxeter number of W.

Remark. Let \(W_{1},\ldots,W_{m}\) be essential finite groups acting in the spaces \(V_{1},\ldots,V_{m}\) and generated by reflections. Let \(C_{j}\) be a chamber relative to \(W_{j}\) . Let W be the group \(W_{1}\times\cdots\times W_{m}\) acting in the space \(V_{1}\times\cdots\times V_{m}\) . Then \(C_{1}\times\cdots\times C_{m}\) is a chamber relative to W. The Coxeter transformations of W defined by C are the products \(c_{1}c_{2}\ldots c_{m}\) , where \(c_{j}\) is a Coxeter transformation of \(W_{j}\) defined by \(C_{j}\) .

